\documentclass[11pt,a4paper]{article}
\usepackage[utf8]{inputenc}
\usepackage[T1]{fontenc}
\usepackage[spanish]{babel}
\usepackage{geometry}
\usepackage{array}
\usepackage{tabularx}
\usepackage{graphicx}
\usepackage{tikz}
\usetikzlibrary{positioning,arrows.meta}
\usepackage{xcolor}
\usepackage{hyperref}
\definecolor{rojoVivo}{RGB}{220,0,0}
\geometry{left=2cm,right=2cm,top=1.8cm,bottom=1.8cm}
\hypersetup{colorlinks=true,urlcolor=rojoVivo,linkcolor=rojoVivo}
\pagestyle{plain}
\setlength{\parindent}{0pt}
\setlength{\parskip}{5pt}
\newcommand{\titulo}[1]{\vspace{7pt}{\large\bfseries\color{rojoVivo}#1}\par\vspace{2pt}}
\newcommand{\encabezado}[1]{\textbf{\color{rojoVivo}#1}}
\begin{document}
\hypersetup{pageanchor=false}
\begin{titlepage}
\centering
\vspace*{0.6cm}
\includegraphics[width=4.2cm]{logo_lasalle_negro.png}\par
\vspace{0.5cm}
{\Large\bfseries Universidad La Salle}\par
\vspace{0.45cm}
{\color{rojoVivo}\rule{0.72\textwidth}{1.4pt}}\par
\vspace{0.65cm}
{\Huge\bfseries\color{rojoVivo} TRABAJO 14}\par
\vspace{0.25cm}
{\Large\bfseries Base de datos de Red Pepper}\par
\vspace{0.25cm}
{\normalsize Identificación y representación de una base de datos}\par
\vspace{0.7cm}
{\bfseries Curso:} Base de Datos I\par
{\bfseries Carrera:} Ingeniería de Software\par
{\bfseries Semestre:} Quinto semestre\par
\vspace{0.65cm}
{\large\bfseries\color{rojoVivo} Autores}\par
\vspace{0.15cm}
Renato Patricio Delgado Huerta\par
Armando Manuel Enriquez Eguiluz\par
Victor Gabriel Lupinta Flores\par
Albana Elizabeth Vasquez Paredes\par
Gabriel Andre Pajuelo Aranibar\par
\vspace{0.55cm}
{\bfseries Docente:} Ricardo Fabrizio Calienes Rodriguez\par
\vfill
2026\par
\end{titlepage}
\hypersetup{pageanchor=true}
\begin{center}
{\LARGE\bfseries\color{rojoVivo} TRABAJO 14}\par
\vspace{3pt}
{\Large\bfseries Base de datos de Red Pepper}\par
\end{center}
\vspace{4pt}

\titulo{1. Descripción de la entidad}
Red Pepper es un local de comida rápida ubicado en la Universidad La Salle, quinto piso. Vende sánguches, empanadas, salchipapas, platos preparados, bandejas para compartir y hamburguesas. Su carta indica atención de lunes a viernes, de 9:00 a.~m. a 5:00 p.~m., y pedidos para recojo.

\titulo{2. Nombre de la base de datos}
\textbf{red\_pepper\_db}. Se propone para registrar clientes, platos y pedidos de manera sencilla.

\titulo{3. Tablas principales}
\begin{tabularx}{\textwidth}{|>{\bfseries}p{2.5cm}|X|p{2.6cm}|}\hline
\encabezado{Tabla} & \encabezado{Qué guarda} & \encabezado{Registros}\\\hline
clientes & Identificación básica de cada cliente & 20\\\hline
platos & Nombre, categoría y precio de cada plato & 30\\\hline
pedidos & Fecha, cliente, plato y cantidad & 20\\\hline
\end{tabularx}

\titulo{4. Estructura de las tablas}
\begin{tabularx}{\textwidth}{|>{\bfseries}p{2.2cm}|X|p{3.3cm}|p{2.6cm}|}\hline
\encabezado{Tabla} & \encabezado{Campo} & \encabezado{Tipo} & \encabezado{Clave}\\\hline
clientes & id\_cliente & INT & PK\\\hline
clientes & nombre & NVARCHAR(100) & ---\\\hline
platos & id\_plato & INT & PK\\\hline
platos & nombre & NVARCHAR(120) & ---\\\hline
platos & categoria & NVARCHAR(40) & ---\\\hline
platos & precio & DECIMAL(10,2) & ---\\\hline
pedidos & id\_pedido & INT & PK\\\hline
pedidos & fecha & DATE & ---\\\hline
pedidos & id\_cliente & INT & FK a clientes\\\hline
pedidos & id\_plato & INT & FK a platos\\\hline
pedidos & cantidad & INT & ---\\\hline
\end{tabularx}

\newpage
\titulo{5. Diagrama entidad--relación}
\begin{center}
\begin{tikzpicture}[
  entidad/.style={draw=rojoVivo,very thick,rectangle,minimum width=4.25cm,align=left,inner sep=8pt},
  rel/.style={draw=rojoVivo,very thick,-{Stealth[length=3mm]}},
  every node/.style={font=\small}
]
\node[entidad] (clientes) at (-4.7,0) {\textcolor{rojoVivo}{\bfseries CLIENTES}\\[3pt]
  id\_cliente (PK)\\nombre};
\node[entidad] (platos) at (4.7,0) {\textcolor{rojoVivo}{\bfseries PLATOS}\\[3pt]
  id\_plato (PK)\\nombre\\categoria\\precio};
\node[entidad] (pedidos) at (0,-4.5) {\textcolor{rojoVivo}{\bfseries PEDIDOS}\\[3pt]
  id\_pedido (PK)\\fecha\\id\_cliente (FK)\\id\_plato (FK)\\cantidad};
\draw[rel] (clientes.south) -- node[left,black]{\bfseries 1 a N} (pedidos.north west);
\draw[rel] (platos.south) -- node[right,black]{\bfseries 1 a N} (pedidos.north east);
\end{tikzpicture}
\end{center}
\vspace{4pt}
Un \textbf{cliente} puede tener varios pedidos. Un \textbf{plato} puede aparecer en varios pedidos. Cada \textbf{pedido} corresponde a un cliente y a un plato. En este modelo básico, cada pedido contiene un solo plato.

\titulo{6. Registro de datos}
El archivo SQL adjunto incluye \textbf{20 clientes, 30 platos y 20 pedidos}. El Excel adjunto presenta los 30 platos con su categoría y precio. Los clientes y pedidos son datos simulados para el ejercicio.

La carta reúne las empanadas por sabores y permite escoger res o pollo en algunas hamburguesas. Estas opciones se registraron por separado para completar los 30 platos sin agregar productos ajenos a la carta.

\titulo{Fuente}
Carta digital de Red Pepper, consultada el 10/11/2026: \url{https://nexlabs.net.pe/Red-Pepper-menu/}
\end{document}
